%==============================
\chapter{Methodology}
\label{cha:methodology}
%==============================

This work investigates the capability of current Large Language Models (LLMs) to convert sequential C code into parallel code.
To this end, we designed a controlled experiment in which a fixed set of sequential programs is submitted to a set of LLMs, and the resulting parallel implementations are evaluated for correctness, performance, and recurring limitations.
This chapter describes the research questions, the programs used as input, the models evaluated, the prompting strategy, the metrics collected, and the execution environment.

%==============================
\section{Research Questions}
\label{sec:rqs}

The central question guiding this study is:

\begin{citacao}
What is the capability of current LLMs in converting sequential C code to parallel C code using OpenMP?
\end{citacao}

To operationalize this question, we define the following secondary research questions (RQs):

\begin{alineas}
\item \textbf{RQ1:} Do the parallel programs generated by the LLMs compile and produce correct results?
\item \textbf{RQ2:} What is the speedup of the generated parallel programs compared to the original sequential programs?
\item \textbf{RQ3:} What are the most common limitations and failures in the generated code—that is, what prevented the generated programs from being correct or efficient?
\item \textbf{RQ4:} How do the different LLMs compare in terms of the correctness and efficiency of the generated parallel code?
\end{alineas}

RQ1 establishes a baseline of functional correctness since a program that compiles but produces incorrect results provides no practical value.
RQ2 addresses performance, which is the core notion of value in parallel computing.
RQ3 provides a qualitative dimension, characterizing the typical failure modes observed.
Finally, RQ4 synthesizes the results of RQ1 and RQ2 across the evaluated models to enable a comparative analysis.

%==============================
\section{Input Programs}
\label{sec:programs}

The sequential programs used as input were drawn from PAMPAR, a parallel benchmark composed of 11 pseudo-applications implemented in C, originally proposed by \citeonline{GARCIA:PAMPAR:2020} to evaluate the performance and energy consumption of different Parallel Programming Interfaces.
PAMPAR was selected because it provides validated sequential implementations together with reference parallel versions written by expert programmers.
The sequential output of each program serves as the ground truth for correctness checking, and the reference parallel version serves as a performance target.

From the 11 available applications, we selected six programs that span a range of data structures, algorithmic complexities, and parallelization challenges.
Table~\ref{tab:programs} summarizes the selected programs.

\begin{table}[ht]
\centering
\caption{Selected input programs from the PAMPAR benchmark.}
\label{tab:programs}
\begin{tabular}{lllll}
\toprule
\textbf{Acronym} & \textbf{Program} & \textbf{Input ($N$)} & \textbf{Complexity} & \textbf{Difficulty} \\
\midrule
DP  & Dot Product                & $10^{10}$  & $O(N)$   & Easy   \\
NI  & Numerical Integration      & $10^{9}$   & $O(N)$   & Easy   \\
MM  & Matrix Multiplication      & 2048       & $O(N^3)$ & Medium \\
OE  & Odd-Even Sort              & $10^{5}$   & $O(N^2)$ & Hard   \\
DFT & Discrete Fourier Transform & $10^{4}$   & $O(N^2)$ & Medium \\
TR  & Turing Ring                & 1280       & $O(N^3)$ & Hard   \\
\bottomrule
\end{tabular}
\fonte{The author.}
\end{table}

DP and NI consist of independent scalar computations that map naturally onto a parallel reduction, making them the most straightforward cases for parallelization.
MM is highly parallelizable but involves more complex memory access patterns across three large matrices; the choice of loop ordering significantly affects cache performance. 
DFT computes a spectrum of $N$ frequency coefficients, each of which is an independent linear combination of the $N$ input samples, making the outer loop trivially parallel.
OE implements the odd-even transposition sort, which alternates between synchronized phases; the parallel implementation must correctly insert barriers between phases.
TR simulates a predator-prey cellular automaton on a 2D toroidal grid using a wavefront pattern: each diagonal of cells is independent of the previous diagonal, but diagonals must be processed in strict order.
This wavefront dependency makes TR the most challenging case for parallelization.

One program from PAMPAR was considered but ultimately excluded: Dijkstra (DJ), which implements the single-source shortest-path algorithm for a dense graph. 
Profiling with the PAMPAR ELAPSEDTIME macro revealed that the sequential file I/O for reading the graph input dominates total execution time, accounting for roughly 99\% of the measured time. 
Under this condition, Amdahl's Law predicts a maximum speedup of approximately 1.01$\times$ regardless of how well the computation itself is parallelized, making DJ unsuitable for evaluating the parallelization capability of LLMs.

%==============================
\section{Evaluated Models}
\label{sec:models}

Five LLMs were evaluated, selected to cover both proprietary frontier models and models with open weights, while remaining accessible at no cost through their official web interfaces.
Table~\ref{tab:models} lists the evaluated models, their configurations, and the date of access.

\begin{table}[ht]
\centering
\caption{Evaluated LLMs and configurations.}
\label{tab:models}
\begin{tabular}{ll}
\toprule
\textbf{Model} & \textbf{Version / Configuration} \\
\midrule
ChatGPT  & GPT-4o (default configuration)                     \\
Gemini   & Gemini 3.1 Pro, Extended Thinking enabled           \\
Claude   & Claude Sonnet 4.6, MAX effort                       \\
Meta AI  & Meta AI with Thinking mode                          \\
DeepSeek & DeepSeek (Instant mode, DeepThink + Search enabled) \\
\bottomrule
\end{tabular}
\fonte{The author.}
\end{table}

All models were accessed through their official web interfaces in June 2026.

%==============================
\section{Prompting Strategy}
\label{sec:prompt}

To ensure that differences in the generated code could be attributed to the models themselves rather than to variations in the prompt, a single prompt template was used and submitted identically to all models.
Each program was submitted in a new, independent conversation to prevent context from previous interactions from influencing the results.
For each model and each program, a single generation attempt was taken as the primary output, reflecting the result the model produces without iterative correction.

The prompt explicitly requests parallelization using OpenMP, which fixes the target technology and makes outputs directly comparable across models:

\begin{citacao}
\textit{I have a sequential C program that I want to parallelize using OpenMP. 
Please convert it to a parallel version while keeping the same logic and producing the same output as the original. 
Return only the complete parallel C code, ready to compile, without explanations.}
\end{citacao}

The full sequential source code of the program was pasted immediately after this instruction in the same message.

%==============================
\section{Metrics and Evaluation Procedure}
\label{sec:metrics}

For each generated program, the following aspects were evaluated:

\begin{alineas}
\item \textbf{Compilation (RQ1):} The generated code was compiled with GCC and OpenMP enabled. Code that failed to compile was recorded as a failure with the error noted.
\item \textbf{Correctness (RQ1):} The output of the parallel program was compared against the output of the original sequential program running on the same input. Matching outputs were considered correct.
\item \textbf{Speedup (RQ2):} execution times were measured and speedup was computed as $S = T_{\text{seq}} / T_{\text{par}}$.
\item \textbf{Hardware counters:} \texttt{perf stat} was used to collect Instructions Per Cycle (IPC) and cache miss rates to support deeper analysis of performance differences.
\item \textbf{Failure characterization (RQ3):} Each case of incorrect or low-performance output was analyzed and classified by cause.
\item \textbf{Model comparison (RQ4):} The results of the previous metrics were tabulated across models to enable a comparative analysis.
\end{alineas}

To obtain stable execution-time measurements, each program was executed 30 times, and the average was used.
Execution times were measured using the ELAPSEDTIME macro provided by PAMPAR, which wraps only the computational region of interest with internal high-resolution timers, excluding program startup and I/O time.
The parallel programs were executed using 11 threads, matching the number of physical cores available minus one core reserved for the operating system (\texttt{taskset -c 1-11}).
All programs were compiled with the flags \texttt{-O2 -fopenmp -DELAPSEDTIME -lm}.

%==============================
\section{Execution Environment}
\label{sec:environment}

All experiments were carried out on a single machine to ensure consistent and comparable measurements.
The hardware and software configuration is described below:

\begin{alineas}
\item \textbf{CPU:} AMD Ryzen 5 5600X, 6 physical cores (12 logical threads with SMT), base clock 3.7~GHz
\item \textbf{Memory:} 32~GB DDR4
\item \textbf{Operating System:} Arch Linux (kernel 7.0.12-arch1-1)
\item \textbf{Compiler:} GCC 16.1.1
\item \textbf{Compilation flags:} \texttt{-O2 -fopenmp -DELAPSEDTIME -lm}
\item \textbf{Thread count:} 11 threads (\texttt{taskset -c 1-11}; core 0 reserved for the OS)
\end{alineas}